\documentclass[aspectratio=169]{beamer}
\usetheme{Copenhagen}
\usecolortheme{whale}
\usefonttheme[onlymath]{serif}

\usepackage{setspace}
\usepackage{xcolor}
\usepackage{graphicx}
\usepackage{multimedia}
\usepackage{fontawesome}
\usepackage[most]{tcolorbox}
\usepackage{empheq}
\usepackage{fancybox}
\usepackage{annotate-equations}
\usepackage{tikz}
\usetikzlibrary{positioning, calc, shapes.geometric, shapes.multipart, 
	shapes, arrows.meta, arrows, 
	decorations.markings, external, trees}
\tikzstyle{Arrow} = [
thick, 
decoration={
	markings,
	mark=at position 1 with {
		\arrow[thick]{latex}
	}
}, 
shorten >= 3pt, preaction = {decorate}
]

\definecolor{teel}{HTML}{00AEB3}
\definecolor{redorange}{HTML}{F26035}
\definecolor{darkred}{HTML}{8B0000}
\definecolor{darkblue}{HTML}{00008B}
\definecolor{darkgreen}{HTML}{06402B}
\definecolor{CarolinaBlue}{HTML}{4B9CD3}
\definecolor{deepsea}{HTML}{123456}
\definecolor{sharedblue}{HTML}{009292}
\definecolor{distalblue}{HTML}{490092}

\usecolortheme[named=CarolinaBlue]{structure}


\title[Revise \& Resubmit]{Responding to Reviewer Comments}
\author[Paul Zivich]{Zivich Research Group}
\date{~}

\setbeamercovered{transparent}
\setbeamertemplate{navigation symbols}{}  % gets rid of the dumb navigation symbols
\setbeamertemplate{page number in head/foot}{\insertframenumber}  % adds slide #
\setbeamertemplate{headline}{}

\AtBeginSection[]{
	\begin{frame}
		\vfill
		\centering
		\begin{beamercolorbox}[sep=8pt,center,shadow=true,rounded=true]{title}
			\usebeamerfont{title}\insertsectionhead\par%
		\end{beamercolorbox}
		\vfill
	\end{frame}
}

\begin{document}

\maketitle

\begin{frame}{Revise and Resubmit (R\&R)}
	Nearly all the peer-review papers go through a `revise and resubmit` phase
	\begin{itemize}
		\item Receive reviewer comments from peers and have an opportunity to respond
	\end{itemize}
	~\\
	This can be a difficult process, especially earlier on\footnote[frame]{I know of someone who says they celebrate when students post a preprint of their work, as the actual publication process can be stochastic and volatile}
	\begin{itemize}
		\item Collection of some advice and practices
	\end{itemize}
\end{frame}

\begin{frame}{Receiving the Decisions}
	You get an email in your inbox that has the decision on a paper and comments. Here is what to do
	\begin{itemize}
		\item Share with your co-authors (or primary mentor) if they are not CC'd on the decision email
		\item Briefly read through the reviewer comments
		\item Unless you can adequately address them all in a couple hours, \textbf{wait at least a day} before revising the paper or drafting responses to the comments\footnote[frame]{Cannot emphasize this enough. Waiting a day will help you adhere to later pieces of advice}
	\end{itemize}
\end{frame}

\begin{frame}{Advice -- General}
	\begin{itemize}
		\item Do \textbf{not} take the comments personally
		\begin{itemize}
			\item This can be discouraging and emotionally draining. It gets easier as you do it more
			\item Reviewer feedback is not a statement of your worth
		\end{itemize}
		\item Wait at least a day before drafting responses
		\begin{itemize}
			\item See previous point for why to take this time
		\end{itemize}
		\item Start by drafting responses to each comment prior to actually editing the manuscript or materials. Here, we are developing a \textit{strategy} for how to respond
		\begin{itemize}
			\item That strategy should be agreed upon by key authors, if not all authors
			\item Once that is agreed upon, then manuscript edits can take place
		\end{itemize}
	\end{itemize}
\end{frame}

\begin{frame}{Advice -- Response Document}
	\begin{itemize}
		\item Copy comments (exactly as they are, including grammatical mistakes) into another document and use a number system for reviewers and their comments
		\item Always put any editor comments (and your responses) at the topic
		\item When finalizing the response document, note the location (line or page number) of that revision. Quoting the revised text may also be reasonable
		\begin{itemize}
			\item Check the journal policies for how they want revisions
		\end{itemize}
	\end{itemize}
\end{frame}

\begin{frame}{Advice -- Writing Responses}
	\begin{itemize}
		\item Be polite. Sometimes reviewers can be cranky,\footnote[frame]{Remember how it feels dealing with a cranky reviewer when you are a peer-reviewer of other's work} but don't return the feeling. You are trying to convince the editor that your paper should be published. Mud-slinging with a reviewer does not help your case
		\item Be responsive. \textit{As a rule of thumb}, you can only provide a push-back of `no' to a single comment per reviewer. Some modification must be made otherwise
		\item Misunderstandings are on the author. If a reviewer does not understand something, then it is on us to make it clearer. In general, a peer-reviewer will read your paper more closely than any other of your future readers
	\end{itemize}
\end{frame}

\begin{frame}{Example}
	\centering
	\includegraphics[width=0.85\linewidth]{images/response_letter1.png}
\end{frame}

\begin{frame}{Example}
	\centering
	\includegraphics[width=0.9\linewidth]{images/response_letter2.png}
\end{frame}

\begin{frame}{Rejections}
	Sometimes you will receive a rejection 
	\begin{itemize}
		\item Before resubmitting, you should consider any peer-review comments you receive
		\begin{itemize}
			\item Do the response document drafting, as if you had gotten an R\&R. Again, follow the wait a day rule
		\end{itemize}
		\item Use that as a guide to whether modifications should be made
		\item Don't take it personal. Editors reject papers for a whole host of different reasons
	\end{itemize}
\end{frame}

\begin{frame}{Open Discussion}
	\centering
	\large
	what is some advice you would give?
\end{frame}	

\end{document}